问题三在问题一的日前购电与储能调度框架上引入了三次新的光伏预报。难点不在于每次都重新求解，而在于区分“已执行的物理状态”和“尚未执行的购电承诺”：后续预报只能改变未执行时段的普通购电量，且调整收益必须足以覆盖下调违约费或上调溢价。若忽略这一不可逆性，模型会把已发生时段当作可回溯决策，从而高估预报更新的价值。

据此，本文采用“日前承诺--信息更新--实时执行”的两层滚动结构。每日 $0{:}00$ 先以当时发布的光伏预报和因果负荷预测制定基准承诺；在 $6{:}00$、$12{:}00$ 和 $18{:}00$，以已观测的储电量为初始状态，仅对剩余时段比较“维持原承诺”和“允许调整”的未来成本。以二者的差额刻画新预报的信息价值，只有该价值超过求解器货币容差时才实施调整。这样既能直接回答“是否需要利用后续预报”，又避免为极小的数值差异频繁改变计划。

光伏预报是逐小时数据，而调度时段为 10 分钟。每次更新先用该时点已观测的实际光伏作为锚点，再与后续整点预报进行因果线性插值，从而保留预报曲线的日内变化，又不读取未来实际光伏。题目未单列负荷预报，故在每日 $0{:}00$ 仅以此前已结束的最近至多四个同星期日逐时段均值构造负荷预测；日内已执行时段由真实观测闭合，未来时段仍使用同一日前预测。该信息口径使计划、更新和实际执行可严格区分。

各次优化的目标、母线平衡和储电量递推均为线性关系，调整费的绝对值可线性化，因此主模型为线性规划。为防止末时段通过耗尽储能虚减当日账单，在目标中加入由下一日因果信息估计的储电量终端价值。求解后从购电承诺锁定、能量平衡、储电量边界、充放电互斥和调整费用五方面核验；最终再以不更新策略和三个单更新时点策略为对照，识别后续预报的总体价值及关键更新时点。
